\documentclass[11pt,a4paper]{article}
\usepackage[utf8]{inputenc}
\usepackage[margin=1in]{geometry}
\usepackage{amsmath,amssymb,amsthm}
\usepackage{listings}
\usepackage{xcolor}
\usepackage{hyperref}

\newtheorem{theorem}{Theorem}

\lstset{
  language=C++,
  basicstyle=\ttfamily\small,
  keywordstyle=\color{blue}\bfseries,
  commentstyle=\color{green!60!black},
  stringstyle=\color{red!70!black},
  numbers=left,
  numberstyle=\tiny\color{gray},
  breaklines=true,
  frame=single,
  tabsize=4,
  showstringspaces=false
}

\title{IOI 2012 -- Crayfish Scrivener}
\author{Solution}
\date{}

\begin{document}
\maketitle

\section{Problem Summary}
Implement a text editor with three operations:
\begin{enumerate}
    \item \textbf{TypeLetter(c)}: Append character $c$ to the current text.
    \item \textbf{UndoCommands(u)}: Undo the last $u$ commands, restoring the text to its state $u$ operations ago. Undo can undo other undos.
    \item \textbf{GetLetter(p)}: Return the $p$-th character (0-indexed) of the current text.
\end{enumerate}
All operations must be efficient ($O(\log N)$ per operation).

\section{Solution}

\subsection{Persistent Linked List with Binary Lifting}
Rather than storing the full string for each version, represent the text as a linked list of character nodes with parent pointers.

\textbf{Version tracking.} Each operation creates a new version. A version stores:
\begin{itemize}
    \item The tail node (last character) of the string, or $-1$ if empty.
    \item The string length.
\end{itemize}

\textbf{TypeLetter(c).} Create a new node with character $c$ and parent = current tail. Set up binary-lifting jump pointers ($\text{jump}[k] = $ ancestor $2^k$ steps back). New version points to this node with length incremented.

\textbf{UndoCommands(u).} The new version copies the version from $u$ steps ago (version index $\text{current} - u$).

\textbf{GetLetter(p).} From the tail, walk back $(\text{len} - 1 - p)$ steps using binary lifting in $O(\log N)$ time. This operation also creates a new version (identical to the current one) since it counts as an operation for undo purposes.

\subsection{Complexity}
\begin{itemize}
    \item \textbf{TypeLetter:} $O(\log N)$ for jump pointers.
    \item \textbf{UndoCommands:} $O(1)$.
    \item \textbf{GetLetter:} $O(\log N)$.
    \item \textbf{Space:} $O(N \log N)$ for jump pointers.
\end{itemize}

\section{Implementation}

\begin{lstlisting}
#include <bits/stdc++.h>
using namespace std;

const int MAXN = 1000005;
const int LOG = 20;

struct Node {
    char ch;
    int jump[LOG];
};

Node nodes[MAXN];
int nodeCount = 0;

struct Version {
    int tail, len;
};

Version versions[MAXN];
int verCount = 0;

void Init() {
    versions[0] = {-1, 0};
    verCount = 1;
}

void TypeLetter(char c) {
    int id = nodeCount++;
    nodes[id].ch = c;
    nodes[id].jump[0] = versions[verCount - 1].tail;
    for (int k = 1; k < LOG; k++) {
        int prev = nodes[id].jump[k - 1];
        nodes[id].jump[k] = (prev == -1) ? -1 : nodes[prev].jump[k - 1];
    }
    versions[verCount] = {id, versions[verCount - 1].len + 1};
    verCount++;
}

void UndoCommands(int u) {
    versions[verCount] = versions[verCount - 1 - u];
    verCount++;
}

char GetLetter(int p) {
    int tail = versions[verCount - 1].tail;
    int len  = versions[verCount - 1].len;
    int steps = len - 1 - p;
    int node = tail;
    for (int k = LOG - 1; k >= 0; k--)
        if (steps >= (1 << k)) {
            node = nodes[node].jump[k];
            steps -= (1 << k);
        }

    // GetLetter counts as an operation for undo
    versions[verCount] = versions[verCount - 1];
    verCount++;

    return nodes[node].ch;
}
\end{lstlisting}

\end{document}
